% Generated by scripts/assemble.py; corrections belong in source/corrections.json.
\examhead{69}{大连理工大学 2024 年数学分析试卷}{69}
% p070-b06; source PDF page 70
\par\Needspace{5\baselineskip}
\blockanchor{p070-b06}
\majorquestion{一}简答题. 每题 6 分, 共 60 分.\par
% p070-b07; source PDF page 70
\blockanchor{p070-b07}
\question{1}用数学语言描述 $\{a_{n}\}$ 不是基本列.\par
% p070-b08; source PDF page 70
\blockanchor{p070-b08}
\question{2}若 $\lim_{x\to 0}\frac{f(x) - f(-x)}{2x} = 0$ ，则 $f^{\prime}(0) = 0$ .此结论是否成立？为什么？\par
% p070-b09; source PDF page 70
\blockanchor{p070-b09}
\question{3}设 $0 < a < b$, 证明不等式 $\frac{2a}{a^2 + b^2} < \frac{\ln b - \ln a}{b - a}$.\par
% p070-b10; source PDF page 70
\blockanchor{p070-b10}
\question{4}已知 $a_{n} = \sqrt[n]{2022^{n} + (-2023)^{n}}$, $n = 1, 2, \cdots$, 求 $\limsup_{n \to \infty} a_{n}$ 和 $\liminf_{n \to \infty} a_{n}$.\par
% p070-b11; source PDF page 70
\blockanchor{p070-b11}
\question{5}$x \to x_0$ 时, $\alpha = o(1)$. 证明: $(1 + \alpha)^{\frac{1}{\alpha}} - e = -\frac{e}{2}\alpha + o(\alpha), x \to x_0$.\par
% p070-b12; source PDF page 70
\blockanchor{p070-b12}
\question{6}证明: $\sin (x^2)$ 在 $\mathbb{R}$ 上不一致连续.\par
% p070-b13; source PDF page 70
\blockanchor{p070-b13}
\question{7}$f(x)$ 在 $\mathbb{R}$ 上连续可微, $h(x,y) = \begin{cases} \frac{f(x) - f(y)}{x - y}, & x \neq y; \\ f'(x), & x = y. \end{cases}$ 证明: $h(x,y)$ 在 $\mathbb{R}^2$ 上连续.\par
% p070-b14; source PDF page 70
\blockanchor{p070-b14}
\question{8}$a, b$ 为正常数, 求由 $\frac{x^2}{a^2} + \frac{y^2}{b^2} = 2x + 18y$ 所围成的面积.\par
% p070-b15; source PDF page 70
\blockanchor{p070-b15}
\question{9}求极限 $\lim_{x\to 0}\left(\frac{1}{x^2} -\frac{\cot x}{x}\right)$\par
% p070-b16; source PDF page 70
\blockanchor{p070-b16}
\question{10}证明: 含参量 $u$ 的反常积分 $\int_{0}^{+\infty} \sqrt{u} e^{-ux^2} \mathrm{d}x$ 在 $u \in (0, +\infty)$ 上不一致收敛.\par
% p070-b17; source PDF page 70
\par\Needspace{5\baselineskip}
\blockanchor{p070-b17}
\majorquestion{二}计算题. 每题 10 分, 共 30 分.\par
% p070-b18; source PDF page 70
\blockanchor{p070-b18}
\question{1}已知 $y = f(x, t)$，其中 $t$ 是由 $F(x, y, t) = 0$ 确定的关于 $x, y$ 的隐函数，$f$ 和 $F$ 有连续的一阶偏导数，求 $\frac{\mathrm{d}y}{\mathrm{d}x}$.\par
% p070-b19; source PDF page 70
\par\Needspace{6\baselineskip}
\blockanchor{p070-b19}
\question{2}设 $C$ 是 $z = \sqrt{a^2 - x^2 - y^2}$ 与 $x = y$ 的交线, 方向由 $\left(\frac{a}{\sqrt{2}}, \frac{a}{\sqrt{2}}, 0\right)$ 到 $\left(-\frac{a}{\sqrt{2}}, -\frac{a}{\sqrt{2}}, 0\right)$. 计算\par
\nopagebreak[4]
% p070-b20; source PDF page 70
\blockanchor{p070-b20}
\displayformula{\int_ {C} \left(z ^ {3} + 3 x ^ {2} y\right) \mathrm{d} x + \left(x ^ {3} + 3 y ^ {2} z\right) \mathrm{d} y + \left(y ^ {3} + 3 z ^ {2} x\right) \mathrm{d} z.}
% p070-b21; source PDF page 70
\blockanchor{p070-b21}
\question{3}求曲线  $\left\{\begin{aligned}x+y-z=2;\\ z^{2}=x^{2}+y^{2}.\end{aligned}\right.$  上距离原点最近的点.\par
% p070-b22; source PDF page 70
\par\Needspace{5\baselineskip}
\blockanchor{p070-b22}
\majorquestion{三}证明题. 每题 12 分, 共 60 分.\par
% p070-b23; source PDF page 70
\blockanchor{p070-b23}
\question{1}证明: $\int_{0}^{+\infty} \frac{x}{1 + x^{6} \sin^{2} x} \, \mathrm{d}x$ 收敛.\par
% p070-b24; source PDF page 70
\par\Needspace{6\baselineskip}
\blockanchor{p070-b24}
\question{2}$f(x)$ 在 $[0,2]$ 上存在三阶连续导数，$f(0) = f(1) = f(2) = 0$。证明：对任意的 $x \in (0,2)$，存在 $c \in (0,2)$，使得\par
\nopagebreak[4]
% p070-b25; source PDF page 70
\blockanchor{p070-b25}
\displayformula{f (x) = \frac {1}{6} x (x - 1) (x - 2) f ^ {\prime \prime \prime} (c).}
% p071-b00; source PDF page 71
\par\Needspace{6\baselineskip}
\blockanchor{p071-b00}
\question{3}$f(x)$ 在 $[a, b]$ 上连续可微, 求证:\par
\nopagebreak[4]
% p071-b01; source PDF page 71
\blockanchor{p071-b01}
\displayformula{\max _ {x \in [ a, b ]} f (x) \leq \frac {1}{b - a} \left| \int_ {a} ^ {b} f (x) \mathrm{d} x \right| + \int_ {a} ^ {b} | f ^ {\prime} (x) | \mathrm{d} x.}
% p071-b02; source PDF page 71
\blockanchor{p071-b02}
\question{4}证明: $f(x) = \sum_{n=0}^{\infty} \frac{\sin(2^n x)}{n!}$ 在 $\mathbb{R}$ 上有任意阶导数, 但不能在 $\mathbb{R}$ 上展开为幂级数.\par
% p071-b03; source PDF page 71
\blockanchor{p071-b03}
\question{5}设 $a_{n} \neq 0, n = 1,2,\dots, \lim_{n\to \infty}n\left(1 - \left|\frac{a_{n + 1}}{a_n}\right|\right) = p > 1$，证明：$\sum_{n=1}^{\infty} a_{n}$ 绝对收敛.\par
